\section*{अध्याय \ref{ch:g3:sorting-living-things} --- सजीवों की छँटाई}

\begin{solution}{exo:g3:sorting-living-things:1}
\omterm{def:g3:sorting-living-things:criterion}{मानदंड} वह नियम है जिससे छँटाई की जाती है। अच्छा: शरीर की कोई ऐसी
ख़ासियत जिसका प्रेक्षण किया जा सके --- ``\omterm{def:g1:animals-around-us:covering}{पंख} हैं'', ``छह \omterm{def:g1:my-body:parts}{टाँगें} हैं''।
जिनसे बचा जाता है: कहाँ रहता है या इस समय क्या कर रहा है --- ``पानी
के पास रहता है'', ``उड़ता है''।
\end{solution}

\begin{solution}{exo:g3:sorting-living-things:2}
\omterm{prop:g3:sorting-living-things:groups}{कीट}: छह \omterm{def:g1:my-body:parts}{टाँगें} और तीन हिस्सों वाला शरीर। \omterm{prop:g3:sorting-living-things:groups}{मछलियाँ}: मीन-पंख और \omterm{def:g1:animals-around-us:covering}{शल्क},
पानी में जीवन। \omterm{prop:g3:sorting-living-things:groups}{पक्षी}: \omterm{def:g1:animals-around-us:covering}{पंख}, डैने और चोंच। \omterm{prop:g3:sorting-living-things:groups}{स्तनधारी}: \omterm{def:g1:animals-around-us:covering}{रोम} या बाल, और ऐसी
माँएँ जो अपने बच्चों को \omterm{def:g2:animals-grow-up:born}{दूध} पिलाती हैं।
\end{solution}

\begin{solution}{exo:g3:sorting-living-things:3}
मधुमक्खी: \omterm{prop:g3:sorting-living-things:groups}{कीट}। सुनहरी मछली: मछली। शुतुरमुर्ग़: \omterm{prop:g3:sorting-living-things:groups}{पक्षी}। गाय: \omterm{prop:g3:sorting-living-things:groups}{स्तनधारी}।
मेंढक: \omterm{prop:g3:sorting-living-things:groups}{उभयचर}।
\end{solution}

\begin{solution}{exo:g3:sorting-living-things:4}
उसके \omterm{def:g1:animals-around-us:covering}{रोम} हैं और वह अपने बच्चों को \omterm{def:g2:animals-grow-up:born}{दूध} पिलाता है, और उसके न \omterm{def:g1:animals-around-us:covering}{पंख} हैं न
चोंच। उड़ना वह करता ज़रूर है, पर उड़ना \omterm{prop:g3:sorting-living-things:groups}{पक्षी} की परिभाषा नहीं है ---
ख़ासियतें कहती हैं: \omterm{prop:g3:sorting-living-things:groups}{स्तनधारी}।
\end{solution}

\begin{solution}{exo:g3:sorting-living-things:5}
उसके \omterm{def:g1:animals-around-us:covering}{शल्क} नहीं हैं और न ही मछली वाला क्लोम-और-मीन-पंख का शरीर --- और
वह अपने बच्चे को \omterm{def:g2:animals-grow-up:born}{दूध} पिलाती है। तैरना उसकी आदत है; उसकी ख़ासियतें
कहती हैं: \omterm{prop:g3:sorting-living-things:groups}{स्तनधारी}।
\end{solution}

\begin{solution}{exo:g3:sorting-living-things:6}
नहीं। \omterm{prop:g3:sorting-living-things:groups}{कीट} का \omterm{def:g3:sorting-living-things:criterion}{मानदंड} छह \omterm{def:g1:my-body:parts}{टाँगें} (और तीन हिस्सों वाला शरीर) है; मकड़ी की
आठ \omterm{def:g1:my-body:parts}{टाँगें} इस पर खरी नहीं उतरतीं। मकड़ी किसी और समूह की है, पहली नज़र
में वह चाहे जैसी लगे।
\end{solution}

\begin{solution}{exo:g3:sorting-living-things:7}
गुलबहार: \omterm{def:g1:plants-around-us:parts}{फूल} और \omterm{def:g2:from-seed-to-plant:seed}{बीज}, मुलायम हरा \omterm{def:g1:plants-around-us:parts}{डंठल}। मॉस: कभी कोई \omterm{def:g1:plants-around-us:parts}{फूल} नहीं, कोई असली
\omterm{def:g1:plants-around-us:parts}{डंठल} नहीं। चीड़: \omterm{def:g1:plants-around-us:tree}{पेड़} (लकड़ी का तना), सुइयों जैसी \omterm{def:g1:plants-around-us:parts}{पत्तियाँ}, \omterm{def:g1:plants-around-us:parts}{फूल} नहीं
पर \omterm{def:g2:from-seed-to-plant:seed}{बीजों} वाले शंकु। गुलाब की झाड़ी: \omterm{def:g1:plants-around-us:parts}{फूल} और \omterm{def:g2:from-seed-to-plant:seed}{बीज}, लकड़ी के \omterm{def:g1:plants-around-us:parts}{डंठल}।
\end{solution}

\begin{solution}{exo:g3:sorting-living-things:8}
उत्तर खुला है। ज़्यादातर गलियों में \omterm{prop:g3:sorting-living-things:groups}{कीट} कहीं आगे जीतते हैं --- चींटियाँ,
मक्खियाँ, मधुमक्खियाँ, भृंग, तितलियाँ --- और दूसरे नंबर पर \omterm{prop:g3:sorting-living-things:groups}{पक्षी}।
\end{solution}

\begin{solution}{exo:g3:sorting-living-things:9}
उसकी सहेली की। ``पालतू बनाम जंगली'' इस आधार पर छाँटता है कि \omterm{def:g1:animals-around-us:animal}{जंतु} लोगों
के साथ कैसे रहते हैं, और यह बदल सकता है --- आवारा बिल्ली घर के भीतर
और बाहर एक ही \omterm{def:g1:animals-around-us:animal}{जंतु} है। ``\omterm{prop:g3:sorting-living-things:groups}{स्तनधारी} बनाम \omterm{prop:g3:sorting-living-things:groups}{पक्षी}'' शरीर की उन ख़ासियतों से
छाँटता है जो जीवन भर टिकी रहती हैं: यही जीवविज्ञानी वाला \omterm{def:g3:sorting-living-things:criterion}{मानदंड} है।
\end{solution}

\begin{solution}{exo:g3:sorting-living-things:10}
चरण 1: आवरण --- \omterm{def:g1:animals-around-us:covering}{पंख} (कसकर जमे हुए), एक चोंच, दो \omterm{def:g1:my-body:parts}{टाँगें}, डैने (मीन-पंख
जैसे)। चरण 2: बच्चे \omterm{def:g2:animals-grow-up:born}{अंडों} से निकलते हैं। चरण 3: ख़ासियतें \omterm{prop:g3:sorting-living-things:groups}{पक्षियों} से
मिलती हैं। चरण 4: ``मछली की तरह तैरता है, उड़ नहीं सकता'' वाले असर
असहमत हैं --- और हार जाते हैं। पेंग्विन \omterm{prop:g3:sorting-living-things:groups}{पक्षी} है।
\end{solution}
